\documentclass[12pt, ukrainian]{article}
\usepackage[a4paper]{geometry}
\setlength\parindent{0mm}
\geometry{top=1.5cm,bottom=1.5cm,left=1.3cm,right=1.5cm}
\pagestyle{empty}
\usepackage[T2A]{fontenc}
\usepackage[utf8]{inputenc}
\usepackage[ukrainian]{babel}
\usepackage{graphicx}
\usepackage{caption}
\usepackage{caption}
\DeclareCaptionFormat{nocaption}{}
\captionsetup{format=nocaption,aboveskip=0pt,belowskip=0pt}
\usepackage[Export]{adjustbox} 
\adjustboxset{max size={0.9\linewidth}{0.9\paperheight}}
\usepackage{verbatim, listings, clrscode3e, fancyvrb}
\def\code#1{\Verb!#1!}
\setlength\parskip{15pt}\begin{document}
%begin
\begin {large}

\textbf{16.11.2024 Контрольна робота, 113, варіант  \No { 1 }}\par


{
%beginitem
1. QoS.

%enditem


\vspace{5mm}
%beginitem
2. Виберіть усі правильні твердження (якщо такі є).\par
\vspace{2mm}

( 1 ) Затримка розповсюдження — це час, необхідний для того, щоб повідомлення було поширене серед усіх вузлів локальної мережі.

( 2 ) Староста збирає заліковки групи, щоб проставити в них печатки в деканаті. Така діяльність старости відповідає задачам транспортного рівня моделі OSI.

( 3 ) Інформація з вікіпедії зазвичай є неавторизованою.

( 4 ) Маршрутизатор зазвичай працює на транспортному рівні.

( 5 ) В аудиторію зайшов співробітник ректорату і представився. Це він виконав авторизацію.

( 6 ) Затримка передачі — це час, який витрачається на відправлення повідомлення в канал зв“язку.

( 7 ) Маршрутизатор зазвичай працює на прикладному рівні.

( 8 ) За інтенсивності трафіка 2 мережа ще може нормально функціонувати, а за інтенсивності 5 та більше варто вживати термінових заходів для підтримки працездатності мережі.

( 9 ) Протокол IP забезпечує надійну передачу даних.

( 10 ) Пересічний користувач під час веб-сьорфінгу породжує потоковий трафік.

( 11 ) Трафік відеоконференції є еластичним.

( 12 ) Староста збирає заліковки групи, щоб проставити в них печатки в деканаті. Така діяльність старости відповідає задачам мережевого рівня моделі OSI.

( 13 ) Протокол TCP має механізм контролю перенавантаження (congestion).

( 14 ) Маршрутизатор зазвичай працює на канальному (data link) рівні.

( 15 ) Затримка очікування – це час, який користувач очікує надходження повідомлення.
%enditem
}
%end
\end {large}

\vspace{4mm}
%end
\newpage
%begin
\begin {large}

\textbf{16.11.2024 Контрольна робота, 113, варіант  \No { 2 }}\par


{
%beginitem
1. Характеристики інформаційних мереж.

%enditem


\vspace{5mm}
%beginitem
2. Виберіть усі правильні твердження (якщо такі є).\par
\vspace{2mm}

( 1 ) Староста збирає заліковки групи, щоб проставити в них печатки в деканаті. Така діяльність старости відповідає задачам прикладного рівня моделі OSI.

( 2 ) Звернення до ChatGPT на іспиті є використанням неавторизованого джерела інформації.

( 3 ) Викладач дозволив на іспиті студентам використовувати тільки власні рукописні конспекти. Це можна віднести до аутентифікації.

( 4 ) Співробітник ректорату на іспиті вирішив перевірити документи студентів. Це він автентифікує студентів.

( 5 ) Коефіцієнт пульсації знаходиться в межах від 0 (не включаючи) до 1.

( 6 ) Передача документів породжує елестичний трафік.

( 7 ) Модель OSI стосується мережі комутації пакетів.

( 8 ) Цифрові сертифікати можуть використовуватись як для аутентифікації, так і для авторизації.

( 9 ) В аудиторію зайшов співробітник ректорату і представився. Це він виконав ідентифікацію.

( 10 ) На фізичному, канальному та мережевому рівнях об“єкти взаємодіють через інтерфейси, а на транспортному та прикладному рівнях через протоколи.

( 11 ) Для передачі документів мережею велика варіація затримки некритична.

( 12 ) Протокол UDP забезпечує надійну передачу даних.

( 13 ) ідентифікує студентів.

( 14 ) Мережевий протокол містить інформацію про кількість даних, передану мережею за певну одиницю часу.

( 15 ) Мережевий протокол є логічним інтерфейсом між різними рівнями стеку протоколів.
%enditem
}
%end
\end {large}

\vspace{4mm}
%end
\newpage
%begin
\begin {large}

\textbf{16.11.2024 Контрольна робота, 113, варіант  \No { 3 }}\par


{
%beginitem
1. Стек протоколів TCP/IP. Співвідношення з моделлю OSI.

%enditem


\vspace{5mm}
%beginitem
2. Виберіть усі правильні твердження (якщо такі є).\par
\vspace{2mm}

( 1 ) Засоби відеоконференцій зазвичай породжують пульсуючий трафік.

( 2 ) Алгоритм дирявого відра це про маршрутизацію пакетів.

( 3 ) В аудиторію зайшов співробітник ректорату і представився. Це він виконав аутентифікацію.

( 4 ) Усі протоколи транспортного рівня вимагають встановлення з“єднання.

( 5 ) Пересічний користувач під час веб-сьорфінгу породжує пульсуючий трафік.

( 6 ) Модель OSI стосується мережі комутації каналів.

( 7 ) На прикладному рівні встановлюється з“єднання між мережевими інтерфейсами пристроїв.

( 8 ) Маршрутизатор зазвичай працює на мережевому рівні.

( 9 ) Співробітник ректорату на іспиті вирішив перевірити документи студентів. Це він авторизує студентів.

( 10 ) FTP є протоколом транспортного рівня.

( 11 ) При передачі повідомлення довжиною 100Кб більше даних буде передано за використання протоколу TCP ніж за використання протоколу UDP.

( 12 ) Велика варіація затримки є критичною для передачі мультимедіа в режимі реального часу.

( 13 ) Мережевий протокол є фізичним інтерфейсом між об“єктами одного рівня.

( 14 ) Профілювання трафіку полягає в тім, що згладжується пульсація трафіку.

( 15 ) Інтенсивність трафіка, рівна 1, вже є критичною для функціонування мережі.
%enditem
}
%end
\end {large}

\vspace{4mm}
%end
\newpage
%begin
\begin {large}

\textbf{16.11.2024 Контрольна робота, 113, варіант  \No { 4 }}\par


{
%beginitem
1. Модель OSI. Рівні та їх задачі. Принципи будови моделі.

%enditem


\vspace{5mm}
%beginitem
2. Виберіть усі правильні твердження (якщо такі є).\par
\vspace{2mm}

( 1 ) Звернення до ChatGPT на іспиті є використанням неаутентифікованого джерела інформації.

( 2 ) BitTorrent – це протокол транспортного рівня.

( 3 ) Протокол SMTP має механізм контролю перенавантаження (congestion).

( 4 ) TCP підтримує багатоадресну передачу даних.

( 5 ) На прикладному рівні встановлюється з“єднання між процесами.

( 6 ) Протокол TCP забезпечує надійну передачу даних.

( 7 ) При передачі повідомлення довжиною 100Кб більше даних буде передано за використання протоколу UDP ніж за використання протоколу TCP.

( 8 ) Протокол  HTTP допускає постійне з“єднання.

( 9 ) В угоді SLA зазвичай вказуються параметри трафіку, які без жодних відхилень у бік погіршення якості мають підтримуватись постачальником послуги.

( 10 ) Засоби відеоконференцій зазвичай породжують потоковий трафік.

( 11 ) Протокол  HTTP є протоколом із збереженням стану (stateful).

( 12 ) З“єднання за протоколом  HTTP завжди є непостійним.

( 13 ) Доволі часто пульсацію трафіку можна зменшити за рахунок використання буферів і збільшення затримки.

( 14 ) Мережевий протокол є логічним інтерфейсом між об“єктами одного рівня.

( 15 ) Еластичний трафік є чутливим до затримок.
%enditem
}
%end
\end {large}

\vspace{4mm}
%end
\newpage
%begin
\begin {large}

\textbf{16.11.2024 Контрольна робота, 113, варіант  \No { 5 }}\par


{
%beginitem
1. Методи забезпечення належного рівня QoS. Керування чергами.

%enditem


\vspace{5mm}
%beginitem
2. Виберіть усі правильні твердження (якщо такі є).\par
\vspace{2mm}

( 1 ) Протокол  HTTP є протоколом без збереження стану (stateless).

( 2 ) З“єднання за протоколом  HTTP забезпечує двосторонній обмін (full duplex).

( 3 ) Затримка передачі — це час, який повідомлення проводить у черзі в очікуванні відправлення в канал зв“язку, поки каналом передається інше повідомлення.

( 4 ) Протокол UDP має механізм контролю перенавантаження (congestion).

( 5 ) Використання допомоги товариша на іспиті є неавторизованим доступом до інформації.

( 6 ) Еластичний трафік не є чутливим до затримок.

( 7 ) TCP можна використовувати для широкомовлення.

( 8 ) Алгоритм дирявого відра дозволяє зменшити пульсації трафіку.

( 9 ) Зазвичай угода  SLA допускає, що зазначені параметри трафіку підтримуються постачальником послуги з певним рівнем ймовірності.

( 10 ) Викладач дозволив на іспиті студентам використовувати тільки власні рукописні конспекти. Це можна віднести до авторизації.

( 11 ) Інформація з вікіпедії часто є неаутентифікованою.

( 12 ) Протокол HTTP має механізм контролю перенавантаження (congestion).

( 13 ) Мережевий протокол містить інформацію про відхилення апраметрів мережі від заданих.

( 14 ) Мультимедійний трафік завжди є чутливим до спотворень.

( 15 ) Співробітник ректорату на іспиті вирішив перевірити документи студентів. Це він
%enditem
}
%end
\end {large}

\vspace{4mm}
%end
\newpage
%begin
\begin {large}

\textbf{16.11.2024 Контрольна робота, 113, варіант  \No { 6 }}\par


{
%beginitem
1. Ідентифікація, аутентифікація, авторизація. Технології аутентифікації.

%enditem


\vspace{5mm}
%beginitem
2. Виберіть усі правильні твердження (якщо такі є).\par
\vspace{2mm}

( 1 ) З“єднання за протоколом  HTTP завжди є постійним.

( 2 ) Протокол  HTTP допускає постійне з“єднання.

( 3 ) Співробітник ректорату на іспиті вирішив перевірити документи студентів. Це він

( 4 ) Мережевий протокол містить інформацію про відхилення апраметрів мережі від заданих.

( 5 ) Затримка розповсюдження — це час, необхідний для того, щоб повідомлення було поширене серед усіх вузлів локальної мережі.

( 6 ) FTP є протоколом транспортного рівня.

( 7 ) Протокол IP забезпечує надійну передачу даних.

( 8 ) Звернення до ChatGPT на іспиті є використанням неавторизованого джерела інформації.

( 9 ) На прикладному рівні встановлюється з“єднання між процесами.

( 10 ) Профілювання трафіку полягає в тім, що згладжується пульсація трафіку.

( 11 ) Викладач дозволив на іспиті студентам використовувати тільки власні рукописні конспекти. Це можна віднести до аутентифікації.

( 12 ) Мультимедійний трафік завжди є чутливим до спотворень.

( 13 ) Затримка передачі — це час, який витрачається на відправлення повідомлення в канал зв“язку.

( 14 ) Велика варіація затримки є критичною для передачі мультимедіа в режимі реального часу.

( 15 ) При передачі повідомлення довжиною 100Кб більше даних буде передано за використання протоколу TCP ніж за використання протоколу UDP.
%enditem
}
%end
\end {large}

\vspace{4mm}
%end
\newpage
%begin
\begin {large}

\textbf{16.11.2024 Контрольна робота, 113, варіант  \No { 7 }}\par


{
%beginitem
1. Маршрутизація, її задачі. Алгоритми та протоколи маршрутизації.

%enditem


\vspace{5mm}
%beginitem
2. Виберіть усі правильні твердження (якщо такі є).\par
\vspace{2mm}

( 1 ) Маршрутизатор зазвичай працює на транспортному рівні.

( 2 ) В аудиторію зайшов співробітник ректорату і представився. Це він виконав авторизацію.

( 3 ) Протокол UDP має механізм контролю перенавантаження (congestion).

( 4 ) З“єднання за протоколом  HTTP завжди є постійним.

( 5 ) Передача документів породжує елестичний трафік.

( 6 ) Зазвичай угода  SLA допускає, що зазначені параметри трафіку підтримуються постачальником послуги з певним рівнем ймовірності.

( 7 ) На фізичному, канальному та мережевому рівнях об“єкти взаємодіють через інтерфейси, а на транспортному та прикладному рівнях через протоколи.

( 8 ) При передачі повідомлення довжиною 100Кб більше даних буде передано за використання протоколу UDP ніж за використання протоколу TCP.

( 9 ) Цифрові сертифікати можуть використовуватись як для аутентифікації, так і для авторизації.

( 10 ) Еластичний трафік не є чутливим до затримок.

( 11 ) Мережевий протокол є логічним інтерфейсом між об“єктами одного рівня.

( 12 ) ідентифікує студентів.

( 13 ) Пересічний користувач під час веб-сьорфінгу породжує пульсуючий трафік.

( 14 ) Протокол TCP забезпечує надійну передачу даних.

( 15 ) Староста збирає заліковки групи, щоб проставити в них печатки в деканаті. Така діяльність старости відповідає задачам прикладного рівня моделі OSI.
%enditem
}
%end
\end {large}

\vspace{4mm}
%end
\newpage
%begin
\begin {large}

\textbf{16.11.2024 Контрольна робота, 113, варіант  \No { 8 }}\par


{
%beginitem
1. Ідентифікація, аутентифікація, авторизація. Технології аутентифікації.

%enditem


\vspace{5mm}
%beginitem
2. Виберіть усі правильні твердження (якщо такі є).\par
\vspace{2mm}

( 1 ) Доволі часто пульсацію трафіку можна зменшити за рахунок використання буферів і збільшення затримки.

( 2 ) Староста збирає заліковки групи, щоб проставити в них печатки в деканаті. Така діяльність старости відповідає задачам мережевого рівня моделі OSI.

( 3 ) З“єднання за протоколом  HTTP завжди є непостійним.

( 4 ) В угоді SLA зазвичай вказуються параметри трафіку, які без жодних відхилень у бік погіршення якості мають підтримуватись постачальником послуги.

( 5 ) Усі протоколи транспортного рівня вимагають встановлення з“єднання.

( 6 ) Модель OSI стосується мережі комутації каналів.

( 7 ) Модель OSI стосується мережі комутації пакетів.

( 8 ) Затримка передачі — це час, який повідомлення проводить у черзі в очікуванні відправлення в канал зв“язку, поки каналом передається інше повідомлення.

( 9 ) Трафік відеоконференції є еластичним.

( 10 ) З“єднання за протоколом  HTTP забезпечує двосторонній обмін (full duplex).

( 11 ) Алгоритм дирявого відра це про маршрутизацію пакетів.

( 12 ) Протокол TCP має механізм контролю перенавантаження (congestion).

( 13 ) Староста збирає заліковки групи, щоб проставити в них печатки в деканаті. Така діяльність старости відповідає задачам транспортного рівня моделі OSI.

( 14 ) Пересічний користувач під час веб-сьорфінгу породжує потоковий трафік.

( 15 ) Маршрутизатор зазвичай працює на мережевому рівні.
%enditem
}
%end
\end {large}

\vspace{4mm}
%end
\newpage
%begin
\begin {large}

\textbf{16.11.2024 Контрольна робота, 113, варіант  \No { 9 }}\par


{
%beginitem
1. QoS.

%enditem


\vspace{5mm}
%beginitem
2. Виберіть усі правильні твердження (якщо такі є).\par
\vspace{2mm}

( 1 ) Співробітник ректорату на іспиті вирішив перевірити документи студентів. Це він авторизує студентів.

( 2 ) Інтенсивність трафіка, рівна 1, вже є критичною для функціонування мережі.

( 3 ) В аудиторію зайшов співробітник ректорату і представився. Це він виконав аутентифікацію.

( 4 ) Використання допомоги товариша на іспиті є неавторизованим доступом до інформації.

( 5 ) Алгоритм дирявого відра дозволяє зменшити пульсації трафіку.

( 6 ) За інтенсивності трафіка 2 мережа ще може нормально функціонувати, а за інтенсивності 5 та більше варто вживати термінових заходів для підтримки працездатності мережі.

( 7 ) Засоби відеоконференцій зазвичай породжують пульсуючий трафік.

( 8 ) На прикладному рівні встановлюється з“єднання між мережевими інтерфейсами пристроїв.

( 9 ) Коефіцієнт пульсації знаходиться в межах від 0 (не включаючи) до 1.

( 10 ) Інформація з вікіпедії часто є неаутентифікованою.

( 11 ) Засоби відеоконференцій зазвичай породжують потоковий трафік.

( 12 ) В аудиторію зайшов співробітник ректорату і представився. Це він виконав ідентифікацію.

( 13 ) Маршрутизатор зазвичай працює на прикладному рівні.

( 14 ) Мережевий протокол містить інформацію про кількість даних, передану мережею за певну одиницю часу.

( 15 ) Звернення до ChatGPT на іспиті є використанням неаутентифікованого джерела інформації.
%enditem
}
%end
\end {large}

\vspace{4mm}
%end
\newpage
%begin
\begin {large}

\textbf{16.11.2024 Контрольна робота, 113, варіант  \No { 10 }}\par


{
%beginitem
1. Маршрутизація, її задачі. Алгоритми та протоколи маршрутизації.

%enditem


\vspace{5mm}
%beginitem
2. Виберіть усі правильні твердження (якщо такі є).\par
\vspace{2mm}

( 1 ) Для передачі документів мережею велика варіація затримки некритична.

( 2 ) Співробітник ректорату на іспиті вирішив перевірити документи студентів. Це він автентифікує студентів.

( 3 ) TCP підтримує багатоадресну передачу даних.

( 4 ) Протокол SMTP має механізм контролю перенавантаження (congestion).

( 5 ) TCP можна використовувати для широкомовлення.

( 6 ) BitTorrent – це протокол транспортного рівня.

( 7 ) Викладач дозволив на іспиті студентам використовувати тільки власні рукописні конспекти. Це можна віднести до авторизації.

( 8 ) Протокол  HTTP є протоколом із збереженням стану (stateful).

( 9 ) Затримка очікування – це час, який користувач очікує надходження повідомлення.

( 10 ) Мережевий протокол є фізичним інтерфейсом між об“єктами одного рівня.

( 11 ) Протокол  HTTP є протоколом без збереження стану (stateless).

( 12 ) Мережевий протокол є логічним інтерфейсом між різними рівнями стеку протоколів.

( 13 ) Протокол HTTP має механізм контролю перенавантаження (congestion).

( 14 ) Маршрутизатор зазвичай працює на канальному (data link) рівні.

( 15 ) Еластичний трафік є чутливим до затримок.
%enditem
}
%end
\end {large}

\vspace{4mm}
%end
\newpage
%begin
\begin {large}

\textbf{16.11.2024 Контрольна робота, 113, варіант  \No { 11 }}\par


{
%beginitem
1. Характеристики інформаційних мереж.

%enditem


\vspace{5mm}
%beginitem
2. Виберіть усі правильні твердження (якщо такі є).\par
\vspace{2mm}

( 1 ) Інформація з вікіпедії зазвичай є неавторизованою.

( 2 ) Протокол UDP забезпечує надійну передачу даних.

( 3 ) Протокол SMTP має механізм контролю перенавантаження (congestion).

( 4 ) Протокол  HTTP є протоколом без збереження стану (stateless).

( 5 ) BitTorrent – це протокол транспортного рівня.

( 6 ) Алгоритм дирявого відра це про маршрутизацію пакетів.

( 7 ) Староста збирає заліковки групи, щоб проставити в них печатки в деканаті. Така діяльність старости відповідає задачам мережевого рівня моделі OSI.

( 8 ) Співробітник ректорату на іспиті вирішив перевірити документи студентів. Це він авторизує студентів.

( 9 ) Мережевий протокол є логічним інтерфейсом між різними рівнями стеку протоколів.

( 10 ) Протокол TCP забезпечує надійну передачу даних.

( 11 ) Маршрутизатор зазвичай працює на мережевому рівні.

( 12 ) При передачі повідомлення довжиною 100Кб більше даних буде передано за використання протоколу UDP ніж за використання протоколу TCP.

( 13 ) Для передачі документів мережею велика варіація затримки некритична.

( 14 ) З“єднання за протоколом  HTTP забезпечує двосторонній обмін (full duplex).

( 15 ) З“єднання за протоколом  HTTP завжди є постійним.
%enditem
}
%end
\end {large}

\vspace{4mm}
%end
\newpage
%begin
\begin {large}

\textbf{16.11.2024 Контрольна робота, 113, варіант  \No { 12 }}\par


{
%beginitem
1. Методи забезпечення належного рівня QoS. Керування чергами.

%enditem


\vspace{5mm}
%beginitem
2. Виберіть усі правильні твердження (якщо такі є).\par
\vspace{2mm}

( 1 ) В аудиторію зайшов співробітник ректорату і представився. Це він виконав ідентифікацію.

( 2 ) З“єднання за протоколом  HTTP завжди є непостійним.

( 3 ) Затримка очікування – це час, який користувач очікує надходження повідомлення.

( 4 ) Протокол TCP має механізм контролю перенавантаження (congestion).

( 5 ) Усі протоколи транспортного рівня вимагають встановлення з“єднання.

( 6 ) Мережевий протокол є фізичним інтерфейсом між об“єктами одного рівня.

( 7 ) Мультимедійний трафік завжди є чутливим до спотворень.

( 8 ) Протокол  HTTP є протоколом із збереженням стану (stateful).

( 9 ) Еластичний трафік не є чутливим до затримок.

( 10 ) Модель OSI стосується мережі комутації каналів.

( 11 ) TCP можна використовувати для широкомовлення.

( 12 ) Профілювання трафіку полягає в тім, що згладжується пульсація трафіку.

( 13 ) Інформація з вікіпедії зазвичай є неавторизованою.

( 14 ) Викладач дозволив на іспиті студентам використовувати тільки власні рукописні конспекти. Це можна віднести до авторизації.

( 15 ) Звернення до ChatGPT на іспиті є використанням неаутентифікованого джерела інформації.
%enditem
}
%end
\end {large}

\vspace{4mm}
%end
\newpage
%begin
\begin {large}

\textbf{16.11.2024 Контрольна робота, 113, варіант  \No { 13 }}\par


{
%beginitem
1. Стек протоколів TCP/IP. Співвідношення з моделлю OSI.

%enditem


\vspace{5mm}
%beginitem
2. Виберіть усі правильні твердження (якщо такі є).\par
\vspace{2mm}

( 1 ) FTP є протоколом транспортного рівня.

( 2 ) Староста збирає заліковки групи, щоб проставити в них печатки в деканаті. Така діяльність старости відповідає задачам транспортного рівня моделі OSI.

( 3 ) Маршрутизатор зазвичай працює на прикладному рівні.

( 4 ) Затримка передачі — це час, який витрачається на відправлення повідомлення в канал зв“язку.

( 5 ) Мережевий протокол є логічним інтерфейсом між об“єктами одного рівня.

( 6 ) Інформація з вікіпедії часто є неаутентифікованою.

( 7 ) Засоби відеоконференцій зазвичай породжують потоковий трафік.

( 8 ) Пересічний користувач під час веб-сьорфінгу породжує потоковий трафік.

( 9 ) Інтенсивність трафіка, рівна 1, вже є критичною для функціонування мережі.

( 10 ) Мережевий протокол містить інформацію про відхилення параметрів мережі від заданих.

( 11 ) Мережевий протокол містить інформацію про кількість даних, передану мережею за певну одиницю часу.

( 12 ) Модель OSI стосується мережі комутації пакетів.

( 13 ) Протокол  HTTP допускає постійне з“єднання.

( 14 ) Використання допомоги товариша на іспиті є неавторизованим доступом до інформації.

( 15 ) Маршрутизатор зазвичай працює на транспортному рівні.
%enditem
}
%end
\end {large}

\vspace{4mm}
%end
\newpage
%begin
\begin {large}

\textbf{16.11.2024 Контрольна робота, 113, варіант  \No { 14 }}\par


{
%beginitem
1. Модель OSI. Рівні та їх задачі. Принципи будови моделі.

%enditem


\vspace{5mm}
%beginitem
2. Виберіть усі правильні твердження (якщо такі є).\par
\vspace{2mm}

( 1 ) Протокол IP забезпечує надійну передачу даних.

( 2 ) Затримка передачі — це час, який повідомлення проводить у черзі в очікуванні відправлення в канал зв“язку, поки каналом передається інше повідомлення.

( 3 ) Затримка розповсюдження — це час, необхідний для того, щоб повідомлення було поширене серед усіх вузлів локальної мережі.

( 4 ) Протокол UDP має механізм контролю перенавантаження (congestion).

( 5 ) Коефіцієнт пульсації знаходиться в межах від 0 (не включаючи) до 1.

( 6 ) Звернення до ChatGPT на іспиті є використанням неавторизованого джерела інформації.

( 7 ) На фізичному, канальному та мережевому рівнях об“єкти взаємодіють через інтерфейси, а на транспортному та прикладному рівнях через протоколи.

( 8 ) ідентифікує студентів.

( 9 ) Маршрутизатор зазвичай працює на канальному (data link) рівні.

( 10 ) Викладач дозволив на іспиті студентам використовувати тільки власні рукописні конспекти. Це можна віднести до аутентифікації.

( 11 ) Доволі часто пульсацію трафіку можна зменшити за рахунок використання буферів і збільшення затримки.

( 12 ) TCP підтримує багатоадресну передачу даних.

( 13 ) На прикладному рівні встановлюється з“єднання між мережевими інтерфейсами пристроїв.

( 14 ) В угоді SLA зазвичай вказуються параметри трафіку, які без жодних відхилень у бік погіршення якості мають підтримуватись постачальником послуги.

( 15 ) Зазвичай угода  SLA допускає, що зазначені параметри трафіку підтримуються постачальником послуги з певним рівнем ймовірності.
%enditem
}
%end
\end {large}

\vspace{4mm}
%end
\newpage
%begin
\begin {large}

\textbf{16.11.2024 Контрольна робота, 113, варіант  \No { 15 }}\par


{
%beginitem
1. Стек протоколів TCP/IP. Співвідношення з моделлю OSI.

%enditem


\vspace{5mm}
%beginitem
2. Виберіть усі правильні твердження (якщо такі є).\par
\vspace{2mm}

( 1 ) Протокол UDP забезпечує надійну передачу даних.

( 2 ) Співробітник ректорату на іспиті вирішив перевірити документи студентів. Це він

( 3 ) В аудиторію зайшов співробітник ректорату і представився. Це він виконав авторизацію.

( 4 ) Еластичний трафік є чутливим до затримок.

( 5 ) Пересічний користувач під час веб-сьорфінгу породжує пульсуючий трафік.

( 6 ) Засоби відеоконференцій зазвичай породжують пульсуючий трафік.

( 7 ) Велика варіація затримки є критичною для передачі мультимедіа в режимі реального часу.

( 8 ) Протокол HTTP має механізм контролю перенавантаження (congestion).

( 9 ) На прикладному рівні встановлюється з“єднання між процесами.

( 10 ) Алгоритм дирявого відра дозволяє зменшити пульсації трафіку.

( 11 ) Передача документів породжує елестичний трафік.

( 12 ) За інтенсивності трафіка 2 мережа ще може нормально функціонувати, а за інтенсивності 5 та більше варто вживати термінових заходів для підтримки працездатності мережі.

( 13 ) В аудиторію зайшов співробітник ректорату і представився. Це він виконав аутентифікацію.

( 14 ) Староста збирає заліковки групи, щоб проставити в них печатки в деканаті. Така діяльність старости відповідає задачам прикладного рівня моделі OSI.

( 15 ) Цифрові сертифікати можуть використовуватись як для аутентифікації, так і для авторизації.
%enditem
}
%end
\end {large}

\vspace{4mm}
%end
\newpage
%begin
\begin {large}

\textbf{16.11.2024 Контрольна робота, 113, варіант  \No { 16 }}\par


{
%beginitem
1. Модель OSI. Рівні та їх задачі. Принципи будови моделі.

%enditem


\vspace{5mm}
%beginitem
2. Виберіть усі правильні твердження (якщо такі є).\par
\vspace{2mm}

( 1 ) Трафік відеоконференції є еластичним.

( 2 ) Співробітник ректорату на іспиті вирішив перевірити документи студентів. Це він автентифікує студентів.

( 3 ) При передачі повідомлення довжиною 100Кб більше даних буде передано за використання протоколу TCP ніж за використання протоколу UDP.

( 4 ) Інтенсивність трафіка, рівна 1, вже є критичною для функціонування мережі.

( 5 ) В аудиторію зайшов співробітник ректорату і представився. Це він виконав авторизацію.

( 6 ) Маршрутизатор зазвичай працює на канальному (data link) рівні.

( 7 ) Протокол TCP забезпечує надійну передачу даних.

( 8 ) Коефіцієнт пульсації знаходиться в межах від 0 (не включаючи) до 1.

( 9 ) Засоби відеоконференцій зазвичай породжують потоковий трафік.

( 10 ) В аудиторію зайшов співробітник ректорату і представився. Це він виконав аутентифікацію.

( 11 ) Протокол UDP забезпечує надійну передачу даних.

( 12 ) Використання допомоги товариша на іспиті є неавторизованим доступом до інформації.

( 13 ) Еластичний трафік не є чутливим до затримок.

( 14 ) Викладач дозволив на іспиті студентам використовувати тільки власні рукописні конспекти. Це можна віднести до аутентифікації.

( 15 ) TCP підтримує багатоадресну передачу даних.
%enditem
}
%end
\end {large}

\vspace{4mm}
%end
\newpage
%begin
\begin {large}

\textbf{16.11.2024 Контрольна робота, 113, варіант  \No { 17 }}\par


{
%beginitem
1. Маршрутизація, її задачі. Алгоритми та протоколи маршрутизації.

%enditem


\vspace{2mm}
%beginitem
2. Виберіть усі правильні твердження (якщо такі є).\par
\vspace{2mm}

( 1 ) Мультимедійний трафік завжди є чутливим до спотворень.

( 2 ) Співробітник ректорату на іспиті вирішив перевірити документи студентів. Це він автентифікує студентів.

( 3 ) Співробітник ректорату на іспиті вирішив перевірити документи студентів. Це він авторизує студентів.

( 4 ) Пересічний користувач під час веб-сьорфінгу породжує пульсуючий трафік.

( 5 ) На фізичному, канальному та мережевому рівнях об“єкти взаємодіють через інтерфейси, а на транспортному та прикладному рівнях через протоколи.

( 6 ) Велика варіація затримки є критичною для передачі мультимедіа в режимі реального часу.

( 7 ) Затримка передачі — це час, який повідомлення проводить у черзі в очікуванні відправлення в канал зв“язку, поки каналом передається інше повідомлення.

( 8 ) Протокол  HTTP є протоколом із збереженням стану (stateful).

( 9 ) Зазвичай угода  SLA допускає, що зазначені параметри трафіку підтримуються постачальником послуги з певним рівнем ймовірності.

( 10 ) Модель OSI стосується мережі комутації пакетів.

( 11 ) Староста збирає заліковки групи, щоб проставити в них печатки в деканаті. Така діяльність старости відповідає задачам мережевого рівня моделі OSI.

( 12 ) Доволі часто пульсацію трафіку можна зменшити за рахунок використання буферів і збільшення затримки.

( 13 ) Алгоритм дирявого відра дозволяє зменшити пульсації трафіку.

( 14 ) Інформація з вікіпедії часто є неаутентифікованою.

( 15 ) За інтенсивності трафіка 2 мережа ще може нормально функціонувати, а за інтенсивності 5 та більше варто вживати термінових заходів для підтримки працездатності мережі.
%enditem
}
%end
\end {large}

\vspace{4mm}
%end
\newpage
%begin
\begin {large}

\textbf{16.11.2024 Контрольна робота, 113, варіант  \No { 18 }}\par


{
%beginitem
1. QoS.

%enditem


\vspace{5mm}
%beginitem
2. Виберіть усі правильні твердження (якщо такі є).\par
\vspace{2mm}

( 1 ) Мережевий протокол містить інформацію про відхилення апраметрів мережі від заданих.

( 2 ) Староста збирає заліковки групи, щоб проставити в них печатки в деканаті. Така діяльність старости відповідає задачам прикладного рівня моделі OSI.

( 3 ) При передачі повідомлення довжиною 100Кб більше даних буде передано за використання протоколу TCP ніж за використання протоколу UDP.

( 4 ) Звернення до ChatGPT на іспиті є використанням неаутентифікованого джерела інформації.

( 5 ) Мережевий протокол є логічним інтерфейсом між різними рівнями стеку протоколів.

( 6 ) Староста збирає заліковки групи, щоб проставити в них печатки в деканаті. Така діяльність старости відповідає задачам транспортного рівня моделі OSI.

( 7 ) Викладач дозволив на іспиті студентам використовувати тільки власні рукописні конспекти. Це можна віднести до авторизації.

( 8 ) Звернення до ChatGPT на іспиті є використанням неавторизованого джерела інформації.

( 9 ) Еластичний трафік є чутливим до затримок.

( 10 ) Протокол UDP забезпечує надійну передачу даних.

( 11 ) Алгоритм дирявого відра це про маршрутизацію пакетів.

( 12 ) Мережевий протокол є логічним інтерфейсом між об“єктами одного рівня.

( 13 ) Мережевий протокол є фізичним інтерфейсом між об“єктами одного рівня.

( 14 ) Мережевий протокол містить інформацію про кількість даних, передану мережею за певну одиницю часу.

( 15 ) BitTorrent – це протокол транспортного рівня.
%enditem
}
%end
\end {large}

\vspace{4mm}
%end
\newpage
%begin
\begin {large}

\textbf{16.11.2024 Контрольна робота, 113, варіант  \No { 19 }}\par


{
%beginitem
1. Методи забезпечення належного рівня QoS. Керування чергами.

%enditem


\vspace{5mm}
%beginitem
2. Виберіть усі правильні твердження (якщо такі є).\par
\vspace{2mm}

( 1 ) TCP можна використовувати для широкомовлення.

( 2 ) Маршрутизатор зазвичай працює на транспортному рівні.

( 3 ) Протокол HTTP має механізм контролю перенавантаження (congestion).

( 4 ) Протокол IP забезпечує надійну передачу даних.

( 5 ) Маршрутизатор зазвичай працює на мережевому рівні.

( 6 ) Маршрутизатор зазвичай працює на прикладному рівні.

( 7 ) Протокол  HTTP допускає постійне з“єднання.

( 8 ) Для передачі документів мережею велика варіація затримки некритична.

( 9 ) Протокол  HTTP є протоколом без збереження стану (stateless).

( 10 ) Цифрові сертифікати можуть використовуватись як для аутентифікації, так і для авторизації.

( 11 ) В аудиторію зайшов співробітник ректорату і представився. Це він виконав ідентифікацію.

( 12 ) З“єднання за протоколом  HTTP завжди є постійним.

( 13 ) Затримка передачі — це час, який витрачається на відправлення повідомлення в канал зв“язку.

( 14 ) Усі протоколи транспортного рівня вимагають встановлення з“єднання.

( 15 ) Модель OSI стосується мережі комутації каналів.
%enditem
}
%end
\end {large}

\vspace{4mm}
%end
\newpage
%begin
\begin {large}

\textbf{16.11.2024 Контрольна робота, 113, варіант  \No { 20 }}\par


{
%beginitem
1. Характеристики інформаційних мереж.

%enditem


\vspace{5mm}
%beginitem
2. Виберіть усі правильні твердження (якщо такі є).\par
\vspace{2mm}

( 1 ) Протокол UDP має механізм контролю перенавантаження (congestion).

( 2 ) Трафік відеоконференції є еластичним.

( 3 ) Затримка розповсюдження — це час, необхідний для того, щоб повідомлення було поширене серед усіх вузлів локальної мережі.

( 4 ) Профілювання трафіку полягає в тім, що згладжується пульсація трафіку.

( 5 ) Протокол SMTP має механізм контролю перенавантаження (congestion).

( 6 ) В угоді SLA зазвичай вказуються параметри трафіку, які без жодних відхилень у бік погіршення якості мають підтримуватись постачальником послуги.

( 7 ) Інформація з вікіпедії зазвичай є неавторизованою.

( 8 ) FTP є протоколом транспортного рівня.

( 9 ) Затримка очікування – це час, який користувач очікує надходження повідомлення.

( 10 ) Засоби відеоконференцій зазвичай породжують пульсуючий трафік.

( 11 ) З“єднання за протоколом  HTTP завжди є непостійним.

( 12 ) На прикладному рівні встановлюється з“єднання між мережевими інтерфейсами пристроїв.

( 13 ) Протокол TCP має механізм контролю перенавантаження (congestion).

( 14 ) Співробітник ректорату на іспиті вирішив перевірити документи студентів. Це він

( 15 ) Передача документів породжує елестичний трафік.
%enditem
}
%end
\end {large}

\vspace{4mm}
%end
\newpage
%begin
\begin {large}

\textbf{16.11.2024 Контрольна робота, 113, варіант  \No { 21 }}\par


{
%beginitem
1. Ідентифікація, аутентифікація, авторизація. Технології аутентифікації.

%enditem


\vspace{5mm}
%beginitem
2. Виберіть усі правильні твердження (якщо такі є).\par
\vspace{2mm}

( 1 ) Пересічний користувач під час веб-сьорфінгу породжує потоковий трафік.

( 2 ) З“єднання за протоколом  HTTP забезпечує двосторонній обмін (full duplex).

( 3 ) На прикладному рівні встановлюється з“єднання між процесами.

( 4 ) При передачі повідомлення довжиною 100Кб більше даних буде передано за використання протоколу UDP ніж за використання протоколу TCP.

( 5 ) Протокол UDP забезпечує надійну передачу даних.

( 6 ) Усі протоколи транспортного рівня вимагають встановлення з“єднання.

( 7 ) Звернення до ChatGPT на іспиті є використанням неавторизованого джерела інформації.

( 8 ) В аудиторію зайшов співробітник ректорату і представився. Це він виконав аутентифікацію.

( 9 ) Трафік відеоконференції є еластичним.

( 10 ) Використання допомоги товариша на іспиті є неавторизованим доступом до інформації.

( 11 ) Інформація з вікіпедії зазвичай є неавторизованою.

( 12 ) Протокол UDP має механізм контролю перенавантаження (congestion).

( 13 ) Староста збирає заліковки групи, щоб проставити в них печатки в деканаті. Така діяльність старости відповідає задачам транспортного рівня моделі OSI.

( 14 ) Пересічний користувач під час веб-сьорфінгу породжує потоковий трафік.

( 15 ) Маршрутизатор зазвичай працює на мережевому рівні.
%enditem
}
%end
\end {large}

\vspace{4mm}
%end
\newpage
%begin
\begin {large}

\textbf{16.11.2024 Контрольна робота, 113, варіант  \No { 22 }}\par


{
%beginitem
1. Модель OSI. Рівні та їх задачі. Принципи будови моделі.

%enditem


\vspace{5mm}
%beginitem
2. Виберіть усі правильні твердження (якщо такі є).\par
\vspace{2mm}

( 1 ) Староста збирає заліковки групи, щоб проставити в них печатки в деканаті. Така діяльність старости відповідає задачам мережевого рівня моделі OSI.

( 2 ) Алгоритм дирявого відра це про маршрутизацію пакетів.

( 3 ) З“єднання за протоколом  HTTP забезпечує двосторонній обмін (full duplex).

( 4 ) На фізичному, канальному та мережевому рівнях об“єкти взаємодіють через інтерфейси, а на транспортному та прикладному рівнях через протоколи.

( 5 ) Викладач дозволив на іспиті студентам використовувати тільки власні рукописні конспекти. Це можна віднести до авторизації.

( 6 ) Зазвичай угода  SLA допускає, що зазначені параметри трафіку підтримуються постачальником послуги з певним рівнем ймовірності.

( 7 ) Цифрові сертифікати можуть використовуватись як для аутентифікації, так і для авторизації.

( 8 ) Для передачі документів мережею велика варіація затримки некритична.

( 9 ) Мережевий протокол є фізичним інтерфейсом між об“єктами одного рівня.

( 10 ) Староста збирає заліковки групи, щоб проставити в них печатки в деканаті. Така діяльність старости відповідає задачам прикладного рівня моделі OSI.

( 11 ) Еластичний трафік є чутливим до затримок.

( 12 ) З“єднання за протоколом  HTTP завжди є постійним.

( 13 ) З“єднання за протоколом  HTTP завжди є непостійним.

( 14 ) Маршрутизатор зазвичай працює на транспортному рівні.

( 15 ) Маршрутизатор зазвичай працює на прикладному рівні.
%enditem
}
%end
\end {large}

\vspace{4mm}
%end
\newpage
%begin
\begin {large}

\textbf{16.11.2024 Контрольна робота, 113, варіант  \No { 23 }}\par


{
%beginitem
1. Ідентифікація, аутентифікація, авторизація. Технології аутентифікації.

%enditem


\vspace{5mm}
%beginitem
2. Виберіть усі правильні твердження (якщо такі є).\par
\vspace{2mm}

( 1 ) TCP підтримує багатоадресну передачу даних.

( 2 ) При передачі повідомлення довжиною 100Кб більше даних буде передано за використання протоколу TCP ніж за використання протоколу UDP.

( 3 ) В аудиторію зайшов співробітник ректорату і представився. Це він виконав ідентифікацію.

( 4 ) Затримка розповсюдження — це час, необхідний для того, щоб повідомлення було поширене серед усіх вузлів локальної мережі.

( 5 ) Мережевий протокол містить інформацію про відхилення апраметрів мережі від заданих.

( 6 ) Інтенсивність трафіка, рівна 1, вже є критичною для функціонування мережі.

( 7 ) Співробітник ректорату на іспиті вирішив перевірити документи студентів. Це він

( 8 ) BitTorrent – це протокол транспортного рівня.

( 9 ) TCP можна використовувати для широкомовлення.

( 10 ) Мережевий протокол є логічним інтерфейсом між об“єктами одного рівня.

( 11 ) Співробітник ректорату на іспиті вирішив перевірити документи студентів. Це він авторизує студентів.

( 12 ) Мультимедійний трафік завжди є чутливим до спотворень.

( 13 ) Профілювання трафіку полягає в тім, що згладжується пульсація трафіку.

( 14 ) За інтенсивності трафіка 2 мережа ще може нормально функціонувати, а за інтенсивності 5 та більше варто вживати термінових заходів для підтримки працездатності мережі.

( 15 ) В аудиторію зайшов співробітник ректорату і представився. Це він виконав авторизацію.
%enditem
}
%end
\end {large}

\vspace{4mm}
%end
\newpage
%begin
\begin {large}

\textbf{16.11.2024 Контрольна робота, 113, варіант  \No { 24 }}\par


{
%beginitem
1. QoS.

%enditem


\vspace{5mm}
%beginitem
2. Виберіть усі правильні твердження (якщо такі є).\par
\vspace{2mm}

( 1 ) Модель OSI стосується мережі комутації пакетів.

( 2 ) Модель OSI стосується мережі комутації каналів.

( 3 ) Пересічний користувач під час веб-сьорфінгу породжує пульсуючий трафік.

( 4 ) Викладач дозволив на іспиті студентам використовувати тільки власні рукописні конспекти. Це можна віднести до аутентифікації.

( 5 ) Доволі часто пульсацію трафіку можна зменшити за рахунок використання буферів і збільшення затримки.

( 6 ) Протокол SMTP має механізм контролю перенавантаження (congestion).

( 7 ) В угоді SLA зазвичай вказуються параметри трафіку, які без жодних відхилень у бік погіршення якості мають підтримуватись постачальником послуги.

( 8 ) Звернення до ChatGPT на іспиті є використанням неаутентифікованого джерела інформації.

( 9 ) ідентифікує студентів.

( 10 ) На прикладному рівні встановлюється з“єднання між мережевими інтерфейсами пристроїв.

( 11 ) Мережевий протокол є логічним інтерфейсом між різними рівнями стеку протоколів.

( 12 ) При передачі повідомлення довжиною 100Кб більше даних буде передано за використання протоколу UDP ніж за використання протоколу TCP.

( 13 ) Засоби відеоконференцій зазвичай породжують пульсуючий трафік.

( 14 ) Алгоритм дирявого відра дозволяє зменшити пульсації трафіку.

( 15 ) Засоби відеоконференцій зазвичай породжують потоковий трафік.
%enditem
}
%end
\end {large}

\vspace{4mm}
%end
\newpage
%begin
\begin {large}

\textbf{16.11.2024 Контрольна робота, 113, варіант  \No { 25 }}\par


{
%beginitem
1. Характеристики інформаційних мереж.

%enditem


\vspace{5mm}
%beginitem
2. Виберіть усі правильні твердження (якщо такі є).\par
\vspace{2mm}

( 1 ) На прикладному рівні встановлюється з“єднання між процесами.

( 2 ) Протокол HTTP має механізм контролю перенавантаження (congestion).

( 3 ) Еластичний трафік не є чутливим до затримок.

( 4 ) Маршрутизатор зазвичай працює на канальному (data link) рівні.

( 5 ) Співробітник ректорату на іспиті вирішив перевірити документи студентів. Це він автентифікує студентів.

( 6 ) Велика варіація затримки є критичною для передачі мультимедіа в режимі реального часу.

( 7 ) Коефіцієнт пульсації знаходиться в межах від 0 (не включаючи) до 1.

( 8 ) Протокол  HTTP є протоколом без збереження стану (stateless).

( 9 ) Затримка передачі — це час, який повідомлення проводить у черзі в очікуванні відправлення в канал зв“язку, поки каналом передається інше повідомлення.

( 10 ) Протокол TCP має механізм контролю перенавантаження (congestion).

( 11 ) Протокол  HTTP є протоколом із збереженням стану (stateful).

( 12 ) Передача документів породжує елестичний трафік.

( 13 ) Протокол  HTTP допускає постійне з“єднання.

( 14 ) Інформація з вікіпедії часто є неаутентифікованою.

( 15 ) FTP є протоколом транспортного рівня.
%enditem
}
%end
\end {large}

\vspace{4mm}
%end
\newpage
%begin
\begin {large}

\textbf{16.11.2024 Контрольна робота, 113, варіант  \No { 26 }}\par


{
%beginitem
1. Методи забезпечення належного рівня QoS. Керування чергами.

%enditem


\vspace{5mm}
%beginitem
2. Виберіть усі правильні твердження (якщо такі є).\par
\vspace{2mm}

( 1 ) Затримка очікування – це час, який користувач очікує надходження повідомлення.

( 2 ) Мережевий протокол містить інформацію про кількість даних, передану мережею за певну одиницю часу.

( 3 ) Затримка передачі — це час, який витрачається на відправлення повідомлення в канал зв“язку.

( 4 ) Протокол IP забезпечує надійну передачу даних.

( 5 ) Протокол TCP забезпечує надійну передачу даних.

( 6 ) Інформація з вікіпедії часто є неаутентифікованою.

( 7 ) Мережевий протокол містить інформацію про кількість даних, передану мережею за певну одиницю часу.

( 8 ) Маршрутизатор зазвичай працює на мережевому рівні.

( 9 ) Для передачі документів мережею велика варіація затримки некритична.

( 10 ) Протокол UDP забезпечує надійну передачу даних.

( 11 ) TCP можна використовувати для широкомовлення.

( 12 ) Протокол TCP забезпечує надійну передачу даних.

( 13 ) Цифрові сертифікати можуть використовуватись як для аутентифікації, так і для авторизації.

( 14 ) На прикладному рівні встановлюється з“єднання між процесами.

( 15 ) Зазвичай угода  SLA допускає, що зазначені параметри трафіку підтримуються постачальником послуги з певним рівнем ймовірності.
%enditem
}
%end
\end {large}

\vspace{4mm}
%end
\newpage
%begin
\begin {large}

\textbf{16.11.2024 Контрольна робота, 113, варіант  \No { 27 }}\par


{
%beginitem
1. Маршрутизація, її задачі. Алгоритми та протоколи маршрутизації.

%enditem


\vspace{5mm}
%beginitem
2. Виберіть усі правильні твердження (якщо такі є).\par
\vspace{2mm}

( 1 ) Алгоритм дирявого відра це про маршрутизацію пакетів.

( 2 ) Затримка розповсюдження — це час, необхідний для того, щоб повідомлення було поширене серед усіх вузлів локальної мережі.

( 3 ) Староста збирає заліковки групи, щоб проставити в них печатки в деканаті. Така діяльність старости відповідає задачам транспортного рівня моделі OSI.

( 4 ) Профілювання трафіку полягає в тім, що згладжується пульсація трафіку.

( 5 ) За інтенсивності трафіка 2 мережа ще може нормально функціонувати, а за інтенсивності 5 та більше варто вживати термінових заходів для підтримки працездатності мережі.

( 6 ) Велика варіація затримки є критичною для передачі мультимедіа в режимі реального часу.

( 7 ) Засоби відеоконференцій зазвичай породжують потоковий трафік.

( 8 ) BitTorrent – це протокол транспортного рівня.

( 9 ) Викладач дозволив на іспиті студентам використовувати тільки власні рукописні конспекти. Це можна віднести до авторизації.

( 10 ) Маршрутизатор зазвичай працює на канальному (data link) рівні.

( 11 ) Протокол  HTTP є протоколом із збереженням стану (stateful).

( 12 ) FTP є протоколом транспортного рівня.

( 13 ) При передачі повідомлення довжиною 100Кб більше даних буде передано за використання протоколу TCP ніж за використання протоколу UDP.

( 14 ) Маршрутизатор зазвичай працює на прикладному рівні.

( 15 ) Співробітник ректорату на іспиті вирішив перевірити документи студентів. Це він
%enditem
}
%end
\end {large}

\vspace{4mm}
%end
\newpage
\end{document}

